\section{模型假设}

为使模型可解且与题目条件相符，本文作如下假设，并逐条说明依据与失效影响。

\begin{itemize}[label=\textbullet,leftmargin=2em]
  \item \textbf{假设一（误差集合而非误差分布）。} 测向误差取有界集合，即真实方位角 $\phi$ 与读数 $\theta$ 满足 $\lvert\phi-\theta\rvert\le\varepsilon$，取 $\varepsilon=1^\circ$ 为主模型，另设 $\varepsilon=1.005^\circ$ 作工程敏感性。\textbf{依据：}题面只给出 $[-1^\circ,1^\circ]$ 这一范围，未给出任何概率分布，且明确指出同一地点误差固定。\textbf{失效影响：}若将来获得经检验的噪声分布，可另建统计模型，但当前不以此改写保证性结论；反之，引入独立同分布假设会使可行域被错误地随重复测量次数收缩。

  \item \textbf{假设二（纯角域多边形与物理先验分段命名）。} 问题一所称"多边形定位区域"解释为各正向测角域的交集 $P$；若再与半径 1800 米目标圆域 $\Omega$ 相交，所得集合一般含圆弧边，另记为 $P_{\text{phys}}$。\textbf{依据：}题目问题一明确称"多边形定位区域"，而实际物理约束含圆形边界。\textbf{失效影响：}若混淆二者，所求直径的对象将发生改变。

  \item \textbf{假设三（首次接收蕴含 $R\ge r$）。} 若首次检测返回有效示向度，则真实距离 $r$ 与接收半径 $R$ 满足 $5<r\le R$。\textbf{依据：}能收到有效信号说明检测点在源的接收范围内，此为题意直接推断。\textbf{失效影响：}忽略该关系会过度收缩候选域或错误判断第二点能否接收；该关系不适用于接收半径随时间变化的模型。

  \item \textbf{假设四（圆的保守外包）。} 数值实现中圆盘一律用外接正多边形代替，以保证不排除真实位置。\textbf{依据：}保持凸性与保守性，便于与半平面统一处理。\textbf{失效影响：}计算量下降但评价结果偏保守；内接多边形或仅取网格中心将失去保证。

  \item \textbf{假设五（同点误差不因重复而减小）。} 在同一位置重复检测同一频道只产生相同约束，不按 $1/\sqrt{N}$ 缩小误差界。\textbf{依据：}题面明确同一地点电磁环境干扰固定。\textbf{失效影响：}若违反，将高估重复测量的信息收益。

  \item \textbf{假设六（舍入语义的保守口径）。} 若 $\pm1^\circ$ 是舍入前的误差界、读数再舍入到两位小数，则用 $\varepsilon=1.005^\circ$ 外包。\textbf{依据：}题面与接口舍入语义未完全明确。\textbf{失效影响：}$\varepsilon$ 取小可能丢失真值，取大则增加保守性；本文已对该口径作敏感性检查。

  \item \textbf{假设七（检测点坐标准确、频道唯一标识目标）。} 检测点坐标由模拟器给出且准确；同一频道只对应一个干扰源。\textbf{依据：}题目"每个干扰源的频道互不相同"及接口语义。\textbf{失效影响：}若频道误关联，交会区域可能为空，此时程序应报告诊断而非丢弃"最不方便"的观测。

  \item \textbf{假设八（源静止、接收半径固定）。} 干扰源位置与接收半径在整场任务中不变。\textbf{依据：}题目未提及源的移动或功率变化。\textbf{失效影响：}正负观测组合所生成的半平面约束依赖"同一源 $R$ 固定"，源变动时不再成立。

  \item \textbf{假设九（隐藏真值不参与决策）。} 模拟器中的真实源位置、类型与朝向仅用于离线生成响应与独立核验，任何执行策略只使用已实际执行的观测。\textbf{依据：}题目中"所有数据无法直接读取"。\textbf{失效影响：}若违规使用，所得时间不构成可执行策略的成绩。
\end{itemize}
